\begin{lstlisting}[language=TypeScript, caption={A route table with a guarded route, nested children, a lazy route table and a constant path.}, label={lst:routes}]
export const routes: Routes = [
  { path: '', component: HomeComponent },
  { path: 'settings', component: SettingsComponent,
    canActivate: [requireAuth],
    children: [{ path: 'profile', component: ProfileComponent }] },
  { path: 'admin', loadChildren: () => import('./admin/admin.routes') },
  { path: paths.about, component: AboutComponent },   // constant, not a literal
];
\end{lstlisting}

\begin{lstlisting}[language=AngularHTML, caption={The root template: the header sits next to the outlet.}, label={lst:shell}]
<app-header />                 <!-- shell: rendered on every screen -->
<router-outlet />              <!-- the active route renders here -->

<!-- header.component.html -->
<a routerLink="/settings">Settings</a>
@if (user()) { <a [routerLink]="['/profile', user().name]">Me</a> }
\end{lstlisting}
\par\smallskip{\footnotesize\noindent\emph{Note.} Its links are present on every screen but belong to none of them.}

\begin{lstlisting}[language=TypeScript, caption={Navigation from an injected service.}, label={lst:service}]
@Injectable({ providedIn: 'root' })
export class AuthService {
  private router = inject(Router);
  private session = inject(SessionStore);
  logout(): void {
    this.session.clear();
    this.router.navigate(['/login']);
  }
}
\end{lstlisting}
\par\smallskip{\footnotesize\noindent\emph{Note.} The call site belongs to no component, so its screens can only be found through the components that call \texttt{logout}.}
